\section{Implementation and comparison details}
\label{app:exp-implementation}

\subsection{V4 optimization and inference}
Table~\ref{tab:exp-recipe} summarizes the settings shared by the four
target-specific fits. The pretrained feature extractors comprise ProtT5
residue representations, the frozen SLEEC scorer, ReactionT5v2, Uni-Mol2,
and ChIRo features. Their pretraining is not replaced by target-only training.
The protein tower retains at most 1,022 residue vectors under the recorded
ends/center policy; the EnzymeMap screening pool already restricts proteins
to at most 650 residues. Feature extraction is reused across compatible
experiments, while downstream retrieval layers and target training dictionaries
remain target-specific.

\begin{table}[htbp]
\centering\small
\caption{Fixed V4 recipe used for the primary benchmark tables. These
settings are distinct from the earlier 18-epoch temperature study and the
longer validation-selected public-baseline training runs.}
\label{tab:exp-recipe}
\begin{tabularx}{\linewidth}{l>{\raggedright\arraybackslash}X}
\toprule
Setting & Value\\\midrule
Seed and phase-1 batch size & 42; 512 association rows\\
Phase-1 retained checkpoint & End of epoch 10\\
Phase-1 optimizer & AdamW; learning rate $10^{-4}$; weight decay $10^{-2}$\\
Phase-1 retrieval objective & Anchor-balanced decoupled multi-positive loss; equal direction weights; fixed inverse temperature $\beta=5$\\
Attention regularizers & Entropy and diversity coefficients both 0.01\\
Phase-2 retained checkpoint & Full-graph update 100; F3 frozen\\
Phase-2 optimizer & AdamW; learning rate $10^{-4}$; weight decay $10^{-3}$; gradient norm clipped at 1\\
Phase-2 retrieval objective & Uniform-positive cross-entropy; temperature 0.2\\
Identity penalty / training residual & 10 / 0.2\\
Biological variant & Phase-2 relative loss; margin 0.1; coefficients $(1,1,1)$; divisor 3\\
Inference multipliers & Protein fusion $\nu=3$; residual $c=0.5$\\
Training-reaction dictionary & Score weight $\alpha=0.40$; 32 protein neighbors; protein and reaction kernel temperatures 0.03\\
Ensemble / hubness correction & Neither used\\
\bottomrule
\end{tabularx}
\end{table}

The dictionary and chemistry normalization are fitted from training inputs.
Test associations are not added to the dictionary. Biological supervision
uses annotations attached to training endpoints; the ReactZyme variant can
also use reaction-only Rhea information. This is additional annotation
supervision, not merely a different way of writing the same association loss.
Conversely, the original V4 chemistry representation already contains
cofactor-related features, so ``without added biology'' does not mean that
every biological input has been removed.

The two active phases have different losses. Phase 1 normalizes positive
contributions by anchor and uses a decoupled denominator. Phase 2 uses a
uniform-positive cross-entropy over the complete training graph. The
original global MLNCE loss is used by the separate parent Horizyn comparator;
it is not the V4 phase-1 objective. Phase-2 biological labels cannot modify
frozen residue queries, and their gradients alone do not establish
biochemical interpretations for the learned views.

\subsection{Public-baseline training adaptations}

The four ReactZyme family implementations are imported from the released
source. Their downstream classifiers use hidden dimension 128, output
dimension 64, dropout 0, batch size 1,000, and AdamW with learning rate
$10^{-4}$ and weight decay $5\times10^{-10}$. Training allows at most 200
epochs, with early stopping after 30 epochs without validation improvement;
the retained checkpoint minimizes validation binary cross-entropy. Four
deterministic corruptions per positive are used because the original negative
pair files are absent. Known training positives are excluded from training
negative samples; validation negatives also exclude validation positives.
Positive and negative weights are 1 and 0.25. This matches downstream
association partitions, but is not an exact reproduction of published
classifier scores.

The released contrastive implementation uses opposite label semantics in
its contrastive and BCE branches. We retain that implementation under the
``contrastive'' label and report a separately named
\texttt{contrastive\_corrected} sensitivity that inverts the contrastive
labels. We do not select between them using test results or relabel the
correction as the unchanged published baseline.

The original Horizyn comparator preserves the public 2,048-to-4,096-to-4,096-to-512
reaction tower and 1,024-to-4,096-to-4,096-to-512 protein tower, global MLNCE
with $\beta=10$, normalized retrieval outputs, and original fingerprint
generators. Its participant-set adapter writes the benchmark components on
one side of the reaction representation. Training uses batch size 16,384,
AdamW learning rate $10^{-4}$, weight decay 0.01, and at most 100 epochs.
Checkpoint selection maximizes mean bidirectional validation all-positive
MRR evaluated every five epochs. The selected epochs are 100, 95, and 100
for Reaction-Sim, Enzyme-Sim, and Time. This is a materially larger training
budget than V4's ten retained epochs, so the observed ordering does not
isolate the loss or architecture effect.

ESM2 and SaProt features follow the recorded released extraction conventions;
29,486 proteins without SaProt structural strings use the released ESM2
fallback. MAT and UniMol inputs preserve the supplied molecular participants
and log conformer failures. The UniMol baseline uses UniMol v1, not the
Uni-Mol2 features supplied to V4. EnzGFM rows use frozen 650M residue means
with the cached 1,022-residue ends/center policy and special tokens excluded;
the published extraction convention differs. These rows are therefore
backbone controls. Reused frozen test-input embeddings do not imply reuse of
target-trained V4 retrieval embeddings.

\paragraph{Completion snapshot.}
The captured evidence contains \completedPublicRuns\ completed public-model
runs: 95 of the 96 unmodified official-family configurations, all 24
corrected-contrastive sensitivities, three Horizyn fits, and four EnzGFM
controls. The remaining unmodified configuration is the Time-split
Transformer with SaProt/UniMol-3D. Three CREEP runs and eight other EnzGFM
configurations did not yet have completed test receipts in this snapshot.
Full structural CLIPZyme and EnzymeCAGE comparisons require complete
validated structures/pockets and justified reaction-input adapters;
VenusRXN requires mapped directional inputs. Their missing results are
not replaced with transferred checkpoints or reduced candidate pools.
The supplementary tables are a frozen snapshot, not a live job-status page.
\FloatBarrier

\section{Metric definitions and retrieval diagnostics}
\label{app:exp-metrics}

\subsection{Released screening metric convention}
For a ranked screening query, let $r_i\in\{0,\ldots,N-1\}$ be the
\emph{zero-based} indices of its $n$ positive candidates and $\rho=n/N$.
The released CLIPZyme notebook implementation reproduced in our evaluator is
\begin{align}
 S_a &= \sum_{i=1}^{n}\exp(-a r_i/N),\\
 D_a &= \frac{\rho(1-\exp(-a))}{\exp(a/N)-1},\\
 F_a &= \frac{\rho\sinh(a/2)}{\cosh(a/2)-\cosh(a/2-a\rho)},\\
 \operatorname{BEDROC}_{a}
     &= \frac{S_a F_a}{D_a}
       + \frac{1}{1-\exp(a(1-\rho))},\qquad a\in\{85,20\}.
 \label{eq:exp-bedroc}
\end{align}
This documents the implemented rank origin explicitly; we do not silently
substitute a one-based formula or clip its finite-size output. The published
CLIPZyme point estimates are closely reproduced under this convention.
The enrichment cutoff is $\lfloor fN\rfloor$, while its denominator remains
$fn$, as in Eq.~\eqref{eq:exp-ef}. We preserve NumPy's recorded descending
score sorting convention for screening, independently of the stable
candidate-index convention in the ReactZyme evaluator. Equal-score handling
must not be silently unified across the two benchmark protocols.

\subsection{All-positive and first-positive retrieval}
In addition to Eq.~\eqref{eq:exp-mrr}, the saved ReactZyme scorer returns
\begin{align}
 \operatorname{MRR}_{\mathrm{first}}
  &=\frac{1}{|Q|}\sum_{q\in Q}
    \frac{1}{\min_{p\in P(q)}\operatorname{rank}_q(p)},\\
 \operatorname{Hit}@k
  &=\frac{1}{|Q|}\sum_{q\in Q}
    \mathbf{1}\!\left[\min_{p\in P(q)}\operatorname{rank}_q(p)\le k\right].
\end{align}
We call the latter Hit@$k$ to avoid importing conflicting Top-$k$/Top-$k$-N
terminology from different papers. Duplicate positive edges are removed
before scoring. Query-level normalization prevents a query with many known
catalysts from contributing proportionally more weight. Queries with no
positive in a specified diagnostic stratum are omitted from that stratum's
mean, while its candidate universe remains fixed.

Table~\ref{tab:exp-diagnostics} shows both reciprocal-rank definitions and
the available hit rates for the same saved score matrices. Mean Rank,
precision-normalized Top-$k$, and Hit@20 are not fabricated from these
summaries; they would require additional rank-level aggregation.

\begin{landscape}
\begin{table}[p]
\centering\small
\caption{Diagnostic ReactZyme metrics from the saved primary test summaries.
``MRR'' is the all-positive metric; ``First MRR'' uses only the earliest
positive. Values are fractions. No checkpoint is reselected using this table.}
\label{tab:exp-diagnostics}
\resizebox{\linewidth}{!}{% Generated from evidence_snapshot.json by prepare_tables.py.
\begin{tabular}{lllrrrrr}
\toprule
Model & Split & Direction & MRR & First MRR & Hit@1 & Hit@5 & Hit@10 \\
\midrule
\methodname & Reaction similarity & R$\to$E & 0.400381 & 0.594198 & 0.528497 & 0.655440 & 0.743523 \\
\methodname & Reaction similarity & E$\to$R & 0.534366 & 0.534383 & 0.436274 & 0.647399 & 0.696623 \\
\methodname & Enzyme similarity & R$\to$E & 0.666249 & 0.926589 & 0.890655 & 0.971392 & 0.984107 \\
\methodname & Enzyme similarity & E$\to$R & 0.971355 & 0.971502 & 0.950195 & 0.994848 & 0.997138 \\
\methodname & Time & R$\to$E & 0.538285 & 0.754115 & 0.689066 & 0.832954 & 0.877752 \\
\methodname & Time & E$\to$R & 0.780669 & 0.780883 & 0.674106 & 0.902338 & 0.936874 \\
\vfourBio & Reaction similarity & R$\to$E & 0.400412 & 0.594189 & 0.528497 & 0.655440 & 0.746114 \\
\vfourBio & Reaction similarity & E$\to$R & 0.534777 & 0.534794 & 0.436887 & 0.647808 & 0.696963 \\
\vfourBio & Enzyme similarity & R$\to$E & 0.666234 & 0.926588 & 0.890655 & 0.971392 & 0.984107 \\
\vfourBio & Enzyme similarity & E$\to$R & 0.967373 & 0.967520 & 0.942295 & 0.994733 & 0.997138 \\
\vfourBio & Time & R$\to$E & 0.538317 & 0.754041 & 0.688686 & 0.833333 & 0.878132 \\
\vfourBio & Time & E$\to$R & 0.800614 & 0.800828 & 0.713692 & 0.902501 & 0.937118 \\
\bottomrule
\end{tabular}
}
\end{table}
\end{landscape}

\section{Case1 results across target-trained checkpoints}
\label{app:exp-case-targets}

\begin{table}[htbp]
\centering\small
\caption{Case1 transfer diagnostics for every target-trained model. Unique
and entry counts are primary-paper recovery at rank 25, out of 12 and 15
respectively. AUROC is conditional on 81 workbook-active and 42 non-detect-only
unique sequences. Reaction-Sim is the primary deployment; the other rows
do not replace it after examining outcomes.}
\label{tab:exp-case-targets}
\resizebox{\linewidth}{!}{% Generated from evidence_snapshot.json by prepare_tables.py.
\begin{tabular}{lrrrrrr}
\toprule
Training split & \multicolumn{3}{c}{\methodname} & \multicolumn{3}{c}{\vfourBio} \\
\midrule
 & Unique & Entries & AUROC & Unique & Entries & AUROC \\
\midrule
Reaction similarity & 10 & 10 & 0.612875 & 11 & 10 & 0.613757 \\
Enzyme similarity & 5 & 8 & 0.584656 & 5 & 8 & 0.585538 \\
Time & 5 & 7 & 0.634333 & 5 & 7 & 0.636390 \\
EnzymeMap & 6 & 6 & 0.470606 & 6 & 6 & 0.470018 \\
\bottomrule
\end{tabular}
}
\end{table}

For the EnzymeMap-trained checkpoint, biological supervision leaves the
unique-paper recovery at 6/12 and changes conditional AUROC from 0.470606
to 0.470018. Thus strong EnzymeMap benchmark performance is not sufficient
to establish useful ranking for this particular literature reaction panel.
The lower transfer performance is retained, rather than reporting only the
best target-trained checkpoint.

The primary evidence tier comprises eight paper sources: the UxaE
engineering study of Shin and colleagues; the binding-pocket work of Xia
and colleagues; the thermostable discovery study of Chen and colleagues;
the hyperthermostable enzyme characterization of Chen and colleagues; the
expression/bioconversion study of Zhao and colleagues; the engineering study
of Wang and colleagues; the coordinated engineering study of Zhu and colleagues;
and the mining/characterization study of Guo and colleagues
\citep{caseS01,caseS02,caseS03,caseS04,caseS05,caseS06,caseS16,caseS17}.
The source audit records which evidence came from an abstract, indexed
full-text section, or supporting-information description. It does not imply
that every construct-level assay in the workbook was independently repeated
or fully revalidated. Full sequence and source mappings accompany the local
audit artifacts identified in the package's source manifest.
\FloatBarrier

\section{Scope of the existing ablation evidence}
\label{app:exp-coverage}

\begin{table}[htbp]
\centering\small
\caption{What is controlled by the recorded ablations. A sensitivity analysis
over nonzero coefficients is not a component-removal experiment.}
\label{tab:exp-coverage}
\begin{tabularx}{\linewidth}{p{.25\linewidth}X}
\toprule
Choice & Available evidence and limit\\\midrule
Dictionary weight & Same weights; $\alpha=0.25,0.40,0.50$ at $c=1$ on both benchmarks. Does not remove the dictionary.\\
Residual strength & Same weights; $c=1$ versus $c=0.5$ at $\alpha=0.40$. Does not remove phase 2.\\
Phase-2 biology & No added labels, correct labels, shuffled labels, and three family-removal controls at fixed F3. Single seed; removal lowers total regularization.\\
Inverse temperature & Paired seeds 17/42/73 at 18 epochs on EnzymeMap. Not a three-seed final-V4 experiment or an MLNCE-versus-anchor-balanced comparison.\\
F3 biology & Paired seeds 42/43/44 with weights 0/0.1/1 and unannotated phase 2. Distinct from the phase-2-only variant.\\
Directional encoder & Historical pilots exist; no complete matched final-V4 directional versus invariant comparison on both benchmarks is established by this audit.\\
Learned residue queries & A fixed-budget comparison of four versus one versus no learned views remains necessary.\\
Architecture versus loss & The stronger parent Horizyn result changes inputs, capacity, objective, batch size, and budget. It is not an isolated loss ablation.\\
\bottomrule
\end{tabularx}
\end{table}

The missing matched controls include dictionary removal ($\alpha=0$),
residual removal ($c=0$), both removals together, and a native fusion
multiplier of one with the other final-recipe choices fixed. Loss, direction,
and learned-query contributions require corresponding matched retraining.
Historical best checkpoints cannot be combined into a causal component table.
We report the existing figures as exploratory controlled comparisons, without
inventing the missing runs or treating all variants as independent replications.
\FloatBarrier

\section{Complete recorded public-baseline results}
\label{app:exp-public}

Tables below retain every completed test receipt at \experimentSnapshot.
Abbreviations: \texttt{esm2} is ESM2-650M; \texttt{saprot} uses the recorded
SaProt/ESM2 fallback; \texttt{enzgfm650} is the adapted frozen EnzGFM backbone.
\texttt{mat\_2d}, \texttt{mat\_3d}, \texttt{unimol\_2d}, and
\texttt{unimol\_3d} identify molecular feature sets. Missing feature fields
in the Horizyn adapter row mean its own native fingerprint/ProtT5 inputs,
not missing model inputs. Epochs are validation-selected. All scores below
are all-positive MRR and use the same complete test catalogs as V4.

\begingroup\footnotesize
\setlength{\tabcolsep}{3pt}
% Generated from evidence_snapshot.json by prepare_tables.py.
\begin{longtable}{p{.32\linewidth}llrrr}
\caption{Completed matched-data runs: Reaction similarity. Every recorded configuration is retained; rows are not selected by test performance.}\label{tab:exp-all-reaction-smi}\\
\toprule
Family / adapter & Protein & Reaction & Epoch & R$\to$E & E$\to$R \\
\midrule\endfirsthead
\toprule
Family / adapter & Protein & Reaction & Epoch & R$\to$E & E$\to$R \\
\midrule\endhead
\bottomrule\endfoot
birnn & esm2 & mat\_2d & 29 & 0.151047 & 0.283470 \\
birnn & esm2 & mat\_3d & 20 & 0.141364 & 0.320658 \\
birnn & esm2 & unimol\_2d & 23 & 0.143440 & 0.318308 \\
birnn & esm2 & unimol\_3d & 29 & 0.142052 & 0.313806 \\
birnn & saprot & mat\_2d & 23 & 0.113894 & 0.293993 \\
birnn & saprot & mat\_3d & 26 & 0.113946 & 0.233916 \\
birnn & saprot & unimol\_2d & 26 & 0.125769 & 0.285104 \\
birnn & saprot & unimol\_3d & 22 & 0.113728 & 0.278580 \\
contrastive\_corrected & esm2 & mat\_2d & 26 & 0.063832 & 0.243126 \\
contrastive\_corrected & esm2 & mat\_3d & 38 & 0.060943 & 0.224593 \\
contrastive\_corrected & esm2 & unimol\_2d & 61 & 0.077039 & 0.245657 \\
contrastive\_corrected & esm2 & unimol\_3d & 72 & 0.067378 & 0.242781 \\
contrastive\_corrected & saprot & mat\_2d & 25 & 0.058795 & 0.228902 \\
contrastive\_corrected & saprot & mat\_3d & 25 & 0.050967 & 0.207046 \\
contrastive\_corrected & saprot & unimol\_2d & 34 & 0.057993 & 0.237400 \\
contrastive\_corrected & saprot & unimol\_3d & 33 & 0.057324 & 0.237389 \\
contrastive & esm2 & mat\_2d & 73 & 0.040553 & 0.239766 \\
contrastive & esm2 & mat\_3d & 67 & 0.043296 & 0.238720 \\
contrastive & esm2 & unimol\_2d & 68 & 0.033340 & 0.178929 \\
contrastive & esm2 & unimol\_3d & 73 & 0.019032 & 0.179270 \\
contrastive & saprot & mat\_2d & 85 & 0.033512 & 0.245238 \\
contrastive & saprot & mat\_3d & 74 & 0.040848 & 0.226484 \\
contrastive & saprot & unimol\_2d & 66 & 0.037882 & 0.315695 \\
contrastive & saprot & unimol\_3d & 68 & 0.030569 & 0.260721 \\
horizyn\_participant\_set & -- & -- & 100 & 0.431815 & 0.534799 \\
mlp & esm2 & mat\_2d & 61 & 0.105690 & 0.218712 \\
mlp & esm2 & mat\_3d & 92 & 0.109936 & 0.258306 \\
mlp & esm2 & unimol\_2d & 60 & 0.086496 & 0.214756 \\
mlp & esm2 & unimol\_3d & 68 & 0.100537 & 0.245804 \\
mlp & saprot & mat\_2d & 96 & 0.102422 & 0.254348 \\
mlp & saprot & mat\_3d & 43 & 0.101501 & 0.207088 \\
mlp & saprot & unimol\_2d & 39 & 0.077519 & 0.290615 \\
mlp & saprot & unimol\_3d & 104 & 0.115039 & 0.279126 \\
transformer & enzgfm650 & mat\_2d & 53 & 0.135390 & 0.296464 \\
transformer & enzgfm650 & mat\_3d & 24 & 0.107978 & 0.266878 \\
transformer & esm2 & mat\_2d & 26 & 0.122008 & 0.231432 \\
transformer & esm2 & mat\_3d & 48 & 0.146864 & 0.266343 \\
transformer & esm2 & unimol\_2d & 52 & 0.125993 & 0.270835 \\
transformer & esm2 & unimol\_3d & 28 & 0.100787 & 0.307673 \\
transformer & saprot & mat\_2d & 47 & 0.112577 & 0.248236 \\
transformer & saprot & mat\_3d & 33 & 0.117991 & 0.228077 \\
transformer & saprot & unimol\_2d & 46 & 0.116645 & 0.263947 \\
transformer & saprot & unimol\_3d & 42 & 0.112881 & 0.254783 \\
\end{longtable}

% Generated from evidence_snapshot.json by prepare_tables.py.
\begin{longtable}{p{.32\linewidth}llrrr}
\caption{Completed matched-data runs: Enzyme similarity. Every recorded configuration is retained; rows are not selected by test performance.}\label{tab:exp-all-enzyme-smi}\\
\toprule
Family / adapter & Protein & Reaction & Epoch & R$\to$E & E$\to$R \\
\midrule\endfirsthead
\toprule
Family / adapter & Protein & Reaction & Epoch & R$\to$E & E$\to$R \\
\midrule\endhead
\bottomrule\endfoot
birnn & esm2 & mat\_2d & 36 & 0.273579 & 0.796985 \\
birnn & esm2 & mat\_3d & 32 & 0.283009 & 0.800479 \\
birnn & esm2 & unimol\_2d & 26 & 0.279413 & 0.797344 \\
birnn & esm2 & unimol\_3d & 22 & 0.254641 & 0.789284 \\
birnn & saprot & mat\_2d & 24 & 0.210042 & 0.733203 \\
birnn & saprot & mat\_3d & 20 & 0.188961 & 0.710997 \\
birnn & saprot & unimol\_2d & 26 & 0.212066 & 0.740293 \\
birnn & saprot & unimol\_3d & 20 & 0.199273 & 0.728435 \\
contrastive\_corrected & esm2 & mat\_2d & 29 & 0.104661 & 0.631571 \\
contrastive\_corrected & esm2 & mat\_3d & 25 & 0.112843 & 0.602566 \\
contrastive\_corrected & esm2 & unimol\_2d & 35 & 0.098874 & 0.615604 \\
contrastive\_corrected & esm2 & unimol\_3d & 48 & 0.135219 & 0.705768 \\
contrastive\_corrected & saprot & mat\_2d & 30 & 0.067626 & 0.527562 \\
contrastive\_corrected & saprot & mat\_3d & 43 & 0.088303 & 0.593146 \\
contrastive\_corrected & saprot & unimol\_2d & 23 & 0.067232 & 0.539864 \\
contrastive\_corrected & saprot & unimol\_3d & 32 & 0.072214 & 0.569024 \\
contrastive & esm2 & mat\_2d & 67 & 0.119823 & 0.666660 \\
contrastive & esm2 & mat\_3d & 56 & 0.096680 & 0.620635 \\
contrastive & esm2 & unimol\_2d & 85 & 0.107420 & 0.650696 \\
contrastive & esm2 & unimol\_3d & 61 & 0.065233 & 0.536463 \\
contrastive & saprot & mat\_2d & 56 & 0.075582 & 0.559432 \\
contrastive & saprot & mat\_3d & 86 & 0.084863 & 0.647888 \\
contrastive & saprot & unimol\_2d & 57 & 0.061696 & 0.538967 \\
contrastive & saprot & unimol\_3d & 71 & 0.067325 & 0.586285 \\
horizyn\_participant\_set & -- & -- & 95 & 0.699701 & 0.976990 \\
mlp & esm2 & mat\_2d & 40 & 0.215499 & 0.707968 \\
mlp & esm2 & mat\_3d & 27 & 0.184238 & 0.678472 \\
mlp & esm2 & unimol\_2d & 92 & 0.250146 & 0.764819 \\
mlp & esm2 & unimol\_3d & 64 & 0.228901 & 0.760377 \\
mlp & saprot & mat\_2d & 88 & 0.209825 & 0.737009 \\
mlp & saprot & mat\_3d & 58 & 0.182154 & 0.682060 \\
mlp & saprot & unimol\_2d & 55 & 0.195097 & 0.706072 \\
mlp & saprot & unimol\_3d & 102 & 0.225738 & 0.736062 \\
transformer & enzgfm650 & mat\_2d & 40 & 0.286377 & 0.798907 \\
transformer & esm2 & mat\_2d & 22 & 0.216317 & 0.725992 \\
transformer & esm2 & mat\_3d & 28 & 0.246319 & 0.770046 \\
transformer & esm2 & unimol\_2d & 27 & 0.231870 & 0.777284 \\
transformer & esm2 & unimol\_3d & 49 & 0.280637 & 0.807310 \\
transformer & saprot & mat\_2d & 45 & 0.211748 & 0.738392 \\
transformer & saprot & mat\_3d & 45 & 0.204797 & 0.720958 \\
transformer & saprot & unimol\_2d & 41 & 0.229146 & 0.753961 \\
transformer & saprot & unimol\_3d & 24 & 0.171613 & 0.701101 \\
\end{longtable}

% Generated from evidence_snapshot.json by prepare_tables.py.
\begin{longtable}{p{.32\linewidth}llrrr}
\caption{Completed matched-data runs: Time. Every recorded configuration is retained; rows are not selected by test performance.}\label{tab:exp-all-time}\\
\toprule
Family / adapter & Protein & Reaction & Epoch & R$\to$E & E$\to$R \\
\midrule\endfirsthead
\toprule
Family / adapter & Protein & Reaction & Epoch & R$\to$E & E$\to$R \\
\midrule\endhead
\bottomrule\endfoot
birnn & esm2 & mat\_2d & 28 & 0.127391 & 0.375185 \\
birnn & esm2 & mat\_3d & 18 & 0.125936 & 0.351964 \\
birnn & esm2 & unimol\_2d & 27 & 0.143554 & 0.405404 \\
birnn & esm2 & unimol\_3d & 22 & 0.134341 & 0.357911 \\
birnn & saprot & mat\_2d & 45 & 0.057523 & 0.147505 \\
birnn & saprot & mat\_3d & 22 & 0.053898 & 0.127134 \\
birnn & saprot & unimol\_2d & 19 & 0.042792 & 0.114786 \\
birnn & saprot & unimol\_3d & 31 & 0.056384 & 0.131066 \\
contrastive\_corrected & esm2 & mat\_2d & 22 & 0.041036 & 0.241179 \\
contrastive\_corrected & esm2 & mat\_3d & 27 & 0.038138 & 0.240529 \\
contrastive\_corrected & esm2 & unimol\_2d & 20 & 0.038952 & 0.231652 \\
contrastive\_corrected & esm2 & unimol\_3d & 27 & 0.035515 & 0.241420 \\
contrastive\_corrected & saprot & mat\_2d & 50 & 0.020605 & 0.108698 \\
contrastive\_corrected & saprot & mat\_3d & 54 & 0.019258 & 0.092859 \\
contrastive\_corrected & saprot & unimol\_2d & 41 & 0.024489 & 0.088000 \\
contrastive\_corrected & saprot & unimol\_3d & 46 & 0.025426 & 0.096020 \\
contrastive & esm2 & mat\_2d & 82 & 0.040501 & 0.255020 \\
contrastive & esm2 & mat\_3d & 72 & 0.036686 & 0.275647 \\
contrastive & esm2 & unimol\_2d & 104 & 0.028089 & 0.274003 \\
contrastive & esm2 & unimol\_3d & 61 & 0.017047 & 0.215914 \\
contrastive & saprot & mat\_2d & 61 & 0.015760 & 0.093115 \\
contrastive & saprot & mat\_3d & 57 & 0.014666 & 0.096924 \\
contrastive & saprot & unimol\_2d & 59 & 0.013870 & 0.101722 \\
contrastive & saprot & unimol\_3d & 53 & 0.011945 & 0.086796 \\
horizyn\_participant\_set & -- & -- & 100 & 0.598203 & 0.840621 \\
mlp & esm2 & mat\_2d & 73 & 0.104607 & 0.284913 \\
mlp & esm2 & mat\_3d & 75 & 0.109812 & 0.296219 \\
mlp & esm2 & unimol\_2d & 74 & 0.103276 & 0.299622 \\
mlp & esm2 & unimol\_3d & 55 & 0.108378 & 0.288156 \\
mlp & saprot & mat\_2d & 101 & 0.056410 & 0.141454 \\
mlp & saprot & mat\_3d & 33 & 0.040105 & 0.093413 \\
mlp & saprot & unimol\_2d & 46 & 0.047138 & 0.104551 \\
mlp & saprot & unimol\_3d & 55 & 0.043459 & 0.098773 \\
transformer & enzgfm650 & mat\_2d & 57 & 0.152498 & 0.386741 \\
transformer & esm2 & mat\_2d & 64 & 0.141887 & 0.394288 \\
transformer & esm2 & mat\_3d & 29 & 0.137216 & 0.355277 \\
transformer & esm2 & unimol\_2d & 45 & 0.148942 & 0.398539 \\
transformer & esm2 & unimol\_3d & 29 & 0.122346 & 0.357867 \\
transformer & saprot & mat\_2d & 38 & 0.047058 & 0.111095 \\
transformer & saprot & mat\_3d & 26 & 0.049423 & 0.130129 \\
transformer & saprot & unimol\_2d & 59 & 0.055180 & 0.139714 \\
\end{longtable}

\endgroup
